% !TEX root = task-3-6-2.tex
% Editable task bank. Append new BankTask blocks at the end of the document.
\documentclass[12pt,a4paper]{article}
\usepackage{physics-textbook}
\usepackage{graphicx,needspace}
\usetikzlibrary{decorations.pathmorphing}
\newif\ifShowSolutions
\ifdefined\TutorEdition
  \ShowSolutionstrue
\else
  \ShowSolutionsfalse % Change to \ShowSolutionstrue to display solutions.
\fi
\newcounter{banktask}
\newsavebox{\BankTaskBox}
\newcommand{\BankTaskContent}[3]{%
  \noindent\textbf{\thebanktask.}\quad #2\par
  \ifShowSolutions
    \Needspace{50mm}% Keep the solution heading with its opening text or diagram.
    #3\par
    \ifdefined\TutorEdition
      \ifcsname TutorSource#1\endcsname
        \nopagebreak\smallskip{\small\csname TutorSource#1\endcsname\par}%
      \fi
    \fi
  \fi
}
\newcommand{\BankTask}[3]{%
  \refstepcounter{banktask}\par\medskip
  \sbox{\BankTaskBox}{\begin{minipage}{\linewidth}%
    \BankTaskContent{#1}{#2}{#3}\end{minipage}}%
  \ifdim\dimexpr\ht\BankTaskBox+\dp\BankTaskBox\relax>\dimexpr\textheight-20mm\relax
    \Needspace{35mm}\BankTaskContent{#1}{#2}{#3}%
  \else
    \Needspace{\dimexpr\ht\BankTaskBox+\dp\BankTaskBox\relax}%
    \noindent\usebox{\BankTaskBox}\par
  \fi
  \medskip
}
\begin{document}
\BookSubsection{3.6.2}{Полярные диэлектрики}


% Source sdamgia-31415: https://phys-ege.sdamgia.ru/problem?id=31415
\BankTask{sdamgia-31415}{
Стеклянные незаряженные кубики~$1$ и $2$ и алюминиевые незаряженные кубики~$3$ и $4$ сблизили вплотную и поместили в электрическое поле, напряженность которого направлена горизонтально влево, как показано в левой части рисунка. Затем кубики раздвинули и уже после этого выключили электрическое поле (правая часть рисунка).\par
$1$. В электрическом поле (левая часть рисунка) часть электронов перешла из кубика~$4$ в кубик~$3$.\par
$2$. Кубик $2$, изображенный в правой части рисунка, не заряжен.\par
$3$. В электрическом поле (левая часть рисунка) правая грань кубика~$2$ заряжена положительно.\par
$4$. Кубик $3$, изображенный в правой части рисунка, имеет положительный заряд.\par
$5$. В электрическом поле (левая часть рисунка) разность потенциалов между левой гранью кубика~$3$ и правой гранью кубика~$4$ равна нулю.\par

\begin{center}\begin{tikzpicture}[book diagram,>=stealth]
\draw(0,0)rectangle(1.2,1.2);\draw(.6,0)--(.6,1.2)(0,.6)--(1.2,.6);\foreach\x/\y/\s in{.3/.9/1,.9/.9/2,.3/.3/3,.9/.3/4}\node at(\x,\y){$\s$};\foreach\y in{.3,.5,.7,.9}\draw[->](-.2,\y)--(-1,\y);\node[above]at(-.7,1){$\vec E$};\draw[dashed](1.65,-.2)--(1.65,1.4)(1.8,-.2)--(1.8,1.4);\foreach\x/\y/\s in{2.1/.8/1,3.1/.8/2,2.1/-.2/3,3.1/-.2/4}{\draw(\x,\y)rectangle++(.6,.6);\node at(\x+.3,\y+.3){$\s$};}
\end{tikzpicture}\end{center}
\par
}{
\smallskip\noindent\textbf{Решение}\par
$1$. Неверно. Электроны под действием электрического поля движутся против силовых линий, поэтому переместятся вправо, тогда часть электронов перейдет с кубика~$3$ в кубик~$4$.\par
$2$. Верно. Стекло является диэлектриком, в котором нет свободных носителей заряда. Поэтому кубики из стекла останутся незаряженными.\par
$3$. Неверно. В кубиках из стекла произойдет поляризация, в результате которой правая часть кубика~$2$ будет иметь отрицательный заряд.\par
$4$. Верно. Поскольку электроны в металлах переместятся против силовых линий, то в кубике~$3$ будет недостаток электронов, он будет заряжен положительно.\par
$5$. Верно. В проводнике, помещенном в электростатическое поле, напряженность станет равной $0$, поэтому и разность потенциалов между левой гранью кубика~$3$ и правой гранью кубика~$4$ равна нулю.\par
Ответ: $245$.\par
}

\end{document}
